This chapter describes the runtime streaming pipeline: how events flow from
ingestion through enrichment, scoring, and decision in near real time. It covers
the messaging topology, the software pattern that enables code reuse across
streaming, batch, and API paths, the fault tolerance layer, and the integration
points with PostgreSQL and the nightly batch rebuild.

% ═══════════════════════════════════════════════════════════
\section{Streaming Topology}
% ═══════════════════════════════════════════════════════════

The pipeline is organized as a linear chain of four RedPanda topics, each
consumed by a dedicated microservice that produces to the next:

\begin{figure}[H]
\centering
\begin{tikzpicture}[
    node distance=0.8cm,
    topic/.style={draw, rounded corners, fill=blue!10, minimum width=3cm,
                  minimum height=0.7cm, font=\small\ttfamily, align=center},
    service/.style={draw, fill=orange!15, minimum width=3cm,
                    minimum height=0.7cm, font=\small, align=center},
    arrow/.style={-{Stealth}, thick}
]
\node[topic]   (raw)      {raw-events};
\node[service, below=of raw]      (enr)      {Enrichment};
\node[topic,   below=of enr]      (enriched) {enriched-events};
\node[service, below=of enriched] (scr)      {Scoring};
\node[topic,   below=of scr]      (scored)   {scored-events};
\node[service, below=of scored]   (dec)      {Decision};
\node[topic,   below=of dec]      (decisions){decisions};

\draw[arrow] (raw)      -- (enr);
\draw[arrow] (enr)      -- (enriched);
\draw[arrow] (enriched) -- (scr);
\draw[arrow] (scr)      -- (scored);
\draw[arrow] (scored)   -- (dec);
\draw[arrow] (dec)      -- (decisions);
\end{tikzpicture}
\caption{Streaming topic chain}
\label{fig:topic-chain}
\end{figure}

All four topics share a single partition key: \texttt{client\_id}. This is a
deliberate constraint. The LSTM model treats each client's transaction history
as a temporal sequence across all four operation types. If operations were split
across separate topics (one per operation type), cross-operation ordering would
be lost --- a retrait followed by a versement five minutes later would arrive on
different topics with no guaranteed relative ordering. A single topic partitioned
by \texttt{client\_id} guarantees that all events for a given client arrive at
the Enrichment Service in the order they occurred, preserving the temporal
signal the LSTM depends on.

\subsection{Broker Choice: RedPanda}

RedPanda provides a Kafka-compatible API without the JVM or ZooKeeper
dependency. This matters in a Docker Compose deployment where every service
competes for memory. A Kafka cluster requires a minimum of three containers
(ZooKeeper + broker + schema registry), while RedPanda runs as a single
container with the same producer/consumer API. The streaming wrappers use
\texttt{kafka-python} and are unaware they are communicating with RedPanda
rather than Kafka --- the switch is transparent at the application level.

\subsection{Consumer Groups and Offset Management}

Each service registers a dedicated consumer group
(\texttt{enrichment-service}, \texttt{scoring-service},
\texttt{decision-service}). Consumer groups provide two guarantees that the
pipeline relies on. First, offset tracking: if a service restarts, it resumes
from the last committed offset rather than replaying the entire topic. Second,
horizontal scalability: adding a second instance of the Enrichment Service
automatically splits partitions across the two consumers, doubling throughput
without code changes. In the current deployment, each service runs as a single
consumer, but the architecture supports scaling out by adding replicas.

% ═══════════════════════════════════════════════════════════
\section{Core Extraction Pattern}
% ═══════════════════════════════════════════════════════════

The system supports three execution paths that all require the same computation:

\begin{enumerate}
    \item \textbf{Streaming} --- Kafka consumer loop, one event at a time,
          writes to Redis buffers
    \item \textbf{Batch} --- Airflow DAG iterating over PostgreSQL rows,
          no Redis interaction
    \item \textbf{Simulation API} --- FastAPI endpoint, synchronous
          request-response, read-only Redis access
\end{enumerate}

Duplicating the enrichment, scoring, or decision logic across these three paths
would create a maintenance burden and a drift risk: any feature computation
change would need to be applied in three places. The solution is a consistent
separation between \emph{core logic} and \emph{transport wrappers}.

\subsection{Core Modules}

Each pipeline stage has a core module that encapsulates all domain logic:

\begin{table}[H]
\centering
\caption{Core modules and their responsibilities}
\label{tab:core-modules}
\begin{tabular}{@{}lll@{}}
\toprule
\textbf{Core} & \textbf{File} & \textbf{Public Method} \\
\midrule
EnrichmentCore & \texttt{src/enrichment/core.py} & \texttt{enrich\_event()} \\
ScoringCore    & \texttt{src/scoring/core.py}     & \texttt{score\_event()} \\
DecisionCore   & \texttt{src/decision/core.py}    & \texttt{decide\_event()} \\
\bottomrule
\end{tabular}
\end{table}

Each core receives its dependencies via constructor injection: a Redis client,
model objects, or configuration. The core exposes a single public method
(fa\c{c}ade pattern) and uses private helpers for sub-steps. Critically, the
cores have no knowledge of Kafka, Airflow, or FastAPI --- they operate on
Python dictionaries and return Python dictionaries.

\subsection{Thin Streaming Wrappers}

The Kafka-specific logic lives in \texttt{src/streaming/}, where each wrapper
file does exactly three things: poll a topic, call the core's public method,
and produce the result to the next topic. The wrapper receives its core
instance via dependency injection at startup. The \texttt{\_\_main\_\_} block
is the sole wiring point: it creates the Redis client, loads models, constructs
the core, injects it into the wrapper, and starts the poll loop.

This separation has a concrete payoff. When the batch pipeline was rewritten
(Section~\ref{sec:batch-integration}), it imported \texttt{EnrichmentCore}
and called the same \texttt{\_compute\_features()} method. When the simulation
API was built, it imported all three cores and wired them with a shared Redis
client. In both cases, zero feature computation code was duplicated.

% ═══════════════════════════════════════════════════════════
\section{Dead Letter Queue and Fault Tolerance}
\label{sec:dlq}
% ═══════════════════════════════════════════════════════════

In a streaming pipeline, a single malformed event can block an entire consumer
if the error is not handled. The fault tolerance layer classifies failures into
two categories and applies different strategies to each.

\subsection{Failure Classification}

\paragraph{Poison pills} (\texttt{KeyError}, \texttt{ValueError},
\texttt{JSONDecodeError}) are structurally invalid events that will never
succeed regardless of how many times they are retried. These are routed to
the DLQ immediately with \texttt{retry\_count=0}.

\paragraph{Transient failures} (\texttt{redis.ConnectionError},
\texttt{psycopg2.OperationalError}, network timeouts) are infrastructure
errors that may resolve on their own. The consumer applies exponential
backoff (1s, 2s, 4s) and retries up to 3 times. If all retries fail, the
event is routed to the DLQ with \texttt{retry\_count=3}.

The distinction matters because the two classes require opposite responses.
Retrying a poison pill wastes time and will never succeed. Immediately
discarding a transient failure loses an event that would have succeeded one
second later when Redis reconnected.

\subsection{ResilientConsumer}

The \texttt{ResilientConsumer} class (\texttt{src/streaming/dlq.py}) wraps all
three streaming services. It intercepts exceptions from the core's public
method, classifies them using an \texttt{isinstance} check against a
\texttt{TRANSIENT\_ERRORS} tuple, applies the appropriate strategy, and
produces failed events to a single \texttt{dlq} topic.

Each DLQ message is wrapped in a metadata envelope:

\begin{table}[H]
\centering
\caption{DLQ envelope fields}
\label{tab:dlq-envelope}
\begin{tabular}{@{}lp{7.5cm}@{}}
\toprule
\textbf{Field} & \textbf{Description} \\
\midrule
\texttt{original\_event}  & The raw event that caused the failure \\
\texttt{source\_topic}    & Topic the event was consumed from \\
\texttt{service}          & Service that failed (enrichment, scoring, decision) \\
\texttt{error\_type}      & Exception class name \\
\texttt{error\_message}   & Exception message string \\
\texttt{stacktrace}       & Full Python traceback \\
\texttt{retry\_count}     & Number of retries attempted before routing to DLQ \\
\texttt{max\_retries}     & Configured maximum (3 for transient errors) \\
\texttt{timestamps}       & First failure and last failure timestamps \\
\bottomrule
\end{tabular}
\end{table}

The envelope preserves enough context for an operator to diagnose the failure,
replay the event after fixing the root cause, and distinguish between
infrastructure outages (transient, clustered in time) and data quality issues
(poison pills, scattered across clients).

\subsection{Structured Logging}

The streaming services use structured JSON logging rather than plain
\texttt{print()} statements. A \texttt{JsonFormatter} class
(\texttt{src/streaming/logger.py}) formats every log line as a JSON object
with a fixed schema: \texttt{timestamp} (UTC ISO 8601), \texttt{service},
\texttt{level}, \texttt{message}, and optional contextual fields
(\texttt{event\_id}, \texttt{client\_id}, \texttt{error\_type},
\texttt{retry\_count}, \texttt{score}, \texttt{alert\_tier}).

A \texttt{setup\_logger(service\_name)} factory function creates a named
logger with the JSON formatter attached. Each streaming service calls
this factory once at startup. The structured format enables log aggregation
tools to filter, group, and alert on specific fields without parsing free-text
messages --- for example, filtering all log lines where
\texttt{error\_type = "redis.ConnectionError"} to detect infrastructure
degradation.

% ═══════════════════════════════════════════════════════════
\section{Enrichment Service}
% ═══════════════════════════════════════════════════════════

The Enrichment Service consumes raw events, computes 30 contrast features by
comparing the current event against the client's historical profile, and
produces enriched events for the Scoring Service.

\subsection{Data Sources}

For each incoming event, the service fetches two categories of data from Redis:

\begin{table}[H]
\centering
\caption{Enrichment Service data sources}
\label{tab:enrichment-sources}
\begin{tabular}{@{}llp{6cm}@{}}
\toprule
\textbf{Source} & \textbf{Redis Key} & \textbf{Content} \\
\midrule
Client profile & \texttt{profile:client:\{id\}} & Batch-computed statistics:
    amount means/stds per operation, operation mix distributions, archetype,
    account age \\
Client buffer  & \texttt{buffer:client:\{id\}} & Last 50 raw event summaries
    (timestamp, amount, operation, payload) for velocity features \\
Employee buffer & \texttt{buffer:employee:\{id\}} & Recent events processed
    by this employee for peer comparison features \\
\bottomrule
\end{tabular}
\end{table}

The client profile is rebuilt nightly by the batch pipeline
(Section~\ref{sec:batch-integration}). The buffers are updated per-event by
the Enrichment Service itself after feature computation.

\subsection{Feature Computation}

The 30 enriched features are organized into categories that reflect how each
feature is derived:

\begin{table}[H]
\centering
\caption{Enriched feature categories}
\label{tab:enriched-features-ch5}
\begin{tabular}{@{}lrl@{}}
\toprule
\textbf{Category} & \textbf{Count} & \textbf{Source} \\
\midrule
Raw event (amount, hour, operation one-hot, branch match) & 7 & Event \\
Amount contrast (z-score, ratio to mean, threshold flags) & 6 & Event + Profile \\
Operation contrast (deviation from expected mix)          & 1 & Event + Profile \\
Timing contrast (hour deviation, day-of-month deviation)  & 2 & Event + Profile \\
Counterparty contrast (known vs. new)                     & 1 & Event payload \\
Velocity (24h/7d counts, cumulative amount, duplicates)   & 5 & Redis buffer \\
Employee (transaction count, client pair count)            & 2 & Redis buffer \\
Context (archetype one-hot, account age)                  & 6 & Profile \\
\bottomrule
\end{tabular}
\end{table}

The term \emph{contrast feature} is central to the design. Raw event fields
(amount, hour, operation type) carry no anomaly signal on their own --- a
10{,}000~DT versement is normal for a big business client but extreme for a
student. Every feature is computed relative to the client's own baseline,
transforming absolute values into deviations that are comparable across
archetypes.

\subsection{Amount Statistics Fallback Chain}

When a client has no 30-day transaction history for a given operation type
(common for infrequent operations like cheque deposits), the amount z-score
computation would fail with a division-by-zero. The service applies a three-level
fallback chain: 30-day statistics first, then 90-day statistics, then the
archetype baseline. This ensures every event receives a meaningful contrast
score regardless of the client's transaction density.

\subsection{Sequence Buffer Management}

After computing the 30 features, the Enrichment Service appends the
\emph{unscaled} feature vector to \texttt{sequence:client:\{id\}} in Redis,
capped at 10 entries. The Scoring Service later retrieves this sequence,
scales it using the training-time \texttt{StandardScaler}, and feeds it to
the LSTM. The features are stored unscaled because the scaler is owned by
the Scoring Service --- storing pre-scaled features would create a coupling
between the Enrichment Service and the scaler version, breaking independence
during model retraining.

% ═══════════════════════════════════════════════════════════
\section{Scoring Service}
% ═══════════════════════════════════════════════════════════

The Scoring Service consumes enriched events, runs both anomaly detection
models, and produces a fused anomaly score with optional SHAP explanations.

\subsection{Score Fusion}

The system runs two complementary models on every event:

\begin{itemize}
    \item \textbf{Autoencoder} --- computes a reconstruction error that
          measures how different the current event is from the ``normal''
          patterns learned during training. The error is converted to a
          percentile rank against the training distribution.
    \item \textbf{LSTM} --- evaluates how surprising the current event is
          given the client's recent sequence of transactions. The prediction
          error is similarly converted to a percentile rank.
\end{itemize}

The two scores are combined using a sigmoid-weighted fusion controlled by
account maturity:

\[
w_{\text{lstm}} = \sigma\!\left(\frac{d - 90}{15}\right)
\]

where $d$ is the number of days since the account was opened, and $\sigma$ is
the logistic function. For accounts younger than ${\sim}75$ days, the
Autoencoder dominates ($w_{\text{lstm}} \approx 0$). For mature accounts
($d > 105$), the LSTM dominates ($w_{\text{lstm}} \approx 1$). The midpoint
of 90 days reflects the minimum history needed for the LSTM's sequence model
to produce reliable predictions.

\[
\text{fused\_score} = (1 - w_{\text{lstm}}) \cdot s_{\text{ae}}
                    + w_{\text{lstm}} \cdot s_{\text{lstm}}
\]

\subsection{Cold-Start Guard}

Clients with fewer than 5 events in their sequence buffer present a distinct
problem. Their velocity features (\texttt{tx\_count\_24h},
\texttt{cumulative\_amount\_24h}) sit far outside the training distribution,
which applied a 21-day warm-up filter. Without a guard, these zero-velocity
vectors produce AE reconstruction errors that map to percentiles above 0.99,
generating false BLOCK-tier alerts for every new client.

The Scoring Service detects cold-start conditions by checking the sequence
buffer length in Redis. When the buffer contains fewer than 5 events, it
returns a neutral score of 0.5 with $w_{\text{lstm}} = 0$ and no SHAP
explanations. This creates a three-tier progression:

\begin{table}[H]
\centering
\caption{Scoring behavior by buffer maturity}
\label{tab:cold-start-ch5}
\begin{tabular}{@{}rll@{}}
\toprule
\textbf{Buffer Size} & \textbf{Scoring Mode} & \textbf{Rationale} \\
\midrule
0--4 events  & Neutral (0.5)    & Insufficient data for meaningful scoring \\
5--9 events  & AE only          & Enough for point anomaly, not for sequence \\
10+ events   & Full AE + LSTM   & Both models have sufficient context \\
\bottomrule
\end{tabular}
\end{table}

\subsection{Conditional SHAP}

SHAP explanations are computationally expensive --- KernelExplainer for the
Autoencoder requires multiple model evaluations per event. To avoid wasting
computation on events that will be classified as INFO, the Scoring Service
only generates SHAP explanations when the fused score exceeds the first
threshold ($T_1 = 0.95$). Events below this threshold are overwhelmingly
normal; explanations for them would never reach the supervisor's screen.

% ═══════════════════════════════════════════════════════════
\section{Decision Service}
% ═══════════════════════════════════════════════════════════

The Decision Service consumes scored events and produces the final alert
decision: an alert tier (INFO, REVIEW, ALERT, or BLOCK) with an explanation
combining model scores and regulatory rule triggers.

\subsection{Alert Tier Assignment}

The base tier is derived from the fused anomaly score using three thresholds:

\begin{table}[H]
\centering
\caption{Score-to-tier mapping}
\label{tab:tier-thresholds}
\begin{tabular}{@{}lll@{}}
\toprule
\textbf{Score Range} & \textbf{Tier} & \textbf{Supervisor Action} \\
\midrule
$< 0.95$ (T\textsubscript{1})                          & INFO   & No action required \\
$\geq 0.95$ and $< 0.99$ (T\textsubscript{2})          & REVIEW & Manual review recommended \\
$\geq 0.99$ and $< 0.999$ (T\textsubscript{3})         & ALERT  & Immediate review required \\
$\geq 0.999$                                            & BLOCK  & Transaction held pending review \\
\bottomrule
\end{tabular}
\end{table}

\subsection{Regulatory Rules}

After the score-based tier is assigned, 8 regulatory rules are evaluated. Each
rule can \emph{escalate} the tier but never \emph{de-escalate} it. This is a
deliberate design constraint: if the model assigns ALERT and a rule would
suggest REVIEW, the tier stays at ALERT. Regulatory flags are always logged
for audit purposes regardless of whether they change the tier.

The escalation-only principle exists because rules encode hard regulatory
requirements (BCT reporting thresholds, LAB/FT structuring patterns) that
must never be overridden by a statistical model. A model might assign a low
score to a transaction that happens to hit a regulatory threshold; the rule
ensures the supervisor still sees it.

Selected rules illustrate the pattern:

\begin{itemize}
    \item \textbf{BCT threshold rule} --- any single transaction exceeding
          10{,}000~DT escalates to at least REVIEW, but only when the amount
          z-score $\geq 2$. The z-score guard prevents routine big-business
          transactions (where 10{,}000~DT is normal) from flooding the
          REVIEW queue.
    \item \textbf{Smurfing detection} --- multiple versements in the
          8{,}000--9{,}999~DT band within 48 hours escalates to ALERT,
          targeting deliberate structuring below the reporting threshold.
    \item \textbf{Cheque cap rule} --- any cheque exceeding 30{,}000~DT
          (the cap established by Law n\textsuperscript{o}41-2024, effective
          February 2025) escalates to BLOCK.
\end{itemize}

\subsection{Cold-Start Tier Cap}

Events flagged as \texttt{cold\_start=True} by the Scoring Service still pass
through all regulatory rules --- a genuinely suspicious transaction from a new
client must still trigger the appropriate flags. However, the Decision Service
caps the final tier at REVIEW for cold-start events. This prevents the neutral
score of 0.5 from being escalated to ALERT or BLOCK by rules alone, which would
generate high-priority alerts for clients who simply have not built enough
history yet. The cap is logged in the decision reasons so supervisors understand
why the tier is limited.

% ═══════════════════════════════════════════════════════════
\section{Sink Consumers}
\label{sec:sinks}
% ═══════════════════════════════════════════════════════════

Two sink consumers persist streaming data to PostgreSQL for historical
querying, dashboard display, and the retraining feedback loop.

\subsection{Raw Events Sink}

The \texttt{raw\_events\_sink} consumes from the \texttt{raw-events} topic
and inserts each event into the \texttt{transactions} table. This provides the
permanent record that the batch pipeline uses for nightly profile rebuilds
and that the dashboard uses for transaction history display.

\subsection{Decisions Sink}

The \texttt{decisions\_sink} consumes from the \texttt{decisions} topic and
performs two writes per event:

\begin{enumerate}
    \item \textbf{PostgreSQL} --- inserts into the \texttt{alerts} table with
          the alert tier, fused score, SHAP explanation (stored as JSONB), and
          all triggered rule names. The \texttt{supervisor\_decision} and
          \texttt{decision\_timestamp} fields are initially \texttt{NULL},
          populated later when the supervisor confirms or rejects the alert
          through the dashboard.
    \item \textbf{Redis} --- increments risk history counters for the client
          (\texttt{risk:client:\{id\}}), tracking the count of alerts per tier
          over rolling windows. These counters feed back into the enrichment
          features for subsequent events.
\end{enumerate}

\subsection{Deliberate FK Omission}

The \texttt{alerts} table does not have a foreign key constraint to the
\texttt{transactions} table, despite both containing the same
\texttt{event\_id}. This is intentional. The two sink consumers run in
parallel: one consumes \texttt{raw-events}, the other consumes
\texttt{decisions}. Because the decision pipeline adds latency (enrichment,
scoring, rule evaluation), the decisions sink may attempt to insert an alert
before the raw events sink has inserted the corresponding transaction. A
foreign key constraint would cause intermittent insertion failures depending
on which consumer is faster at any given moment. The \texttt{event\_id} field
still links the two tables logically; the constraint is enforced at the
application level rather than the database level.


% ═══════════════════════════════════════════════════════════
\section{Data Persistence Layer}
\label{sec:persistence}
% ═══════════════════════════════════════════════════════════

The PostgreSQL database serves as the cold store for the entire system. It
receives writes from two streaming sinks and the batch pipeline, and provides
reads for the dashboard, the batch profile rebuild, and the retraining
pipeline. Figure~\ref{fig:er-diagram} shows the table relationships.

\begin{figure}[H]
\centering
\resizebox{\textwidth}{!}{%
\begin{tikzpicture}[
    node distance=1.2cm and 1.8cm,
    every node/.style={font=\small},
    table/.style={
        draw, rectangle, rounded corners=2pt,
        inner sep=0pt, outer sep=0pt,
        minimum width=3.4cm
    },
    header/.style={
        fill=blue!15, font=\small\bfseries\ttfamily,
        minimum width=3.4cm, minimum height=0.5cm, align=center
    },
    cols/.style={
        fill=white, font=\tiny\ttfamily,
        minimum width=3.4cm, align=left, inner sep=4pt
    },
    fk/.style={-{Stealth}, thick},
    nofk/.style={-{Stealth}, thick, dashed},
]

% ── clients_master ──
\node[table] (clients) {
    \begin{tabular}{@{}c@{}}
    \tikz\node[header]{clients\_master};\\[-0.4pt]
    \tikz\node[cols]{%
        \underline{client\_id} : UUID PK\\
        archetype : VARCHAR\\
        account\_opening\_date : DATE\\
        nationality : VARCHAR\\
        client\_type : VARCHAR\\
        home\_branch\_id : VARCHAR
    };
    \end{tabular}
};

% ── employees_master ──
\node[table, right=3cm of clients] (employees) {
    \begin{tabular}{@{}c@{}}
    \tikz\node[header]{employees\_master};\\[-0.4pt]
    \tikz\node[cols]{%
        \underline{employee\_id} : VARCHAR PK\\
        branch\_id : VARCHAR
    };
    \end{tabular}
};

% ── accounts ──
\node[table, below=of clients] (accounts) {
    \begin{tabular}{@{}c@{}}
    \tikz\node[header]{accounts};\\[-0.4pt]
    \tikz\node[cols]{%
        \underline{account\_id} : UUID PK\\
        client\_id : UUID FK\\
        opening\_date : DATE\\
        account\_type : VARCHAR\\
        status : VARCHAR
    };
    \end{tabular}
};

% ── transactions ──
\node[table, right=3cm of accounts] (transactions) {
    \begin{tabular}{@{}c@{}}
    \tikz\node[header]{transactions};\\[-0.4pt]
    \tikz\node[cols]{%
        \underline{event\_id} : UUID PK\\
        client\_id : UUID\\
        account\_id : UUID\\
        employee\_id : VARCHAR\\
        branch\_id : VARCHAR\\
        timestamp : TIMESTAMP\\
        amount : NUMERIC\\
        operation\_type : VARCHAR\\
        payload : JSONB
    };
    \end{tabular}
};

% ── alerts ──
\node[table, below=of accounts] (alerts) {
    \begin{tabular}{@{}c@{}}
    \tikz\node[header]{alerts};\\[-0.4pt]
    \tikz\node[cols]{%
        \underline{alert\_id} : SERIAL PK\\
        event\_id : UUID\\
        client\_id : UUID\\
        tier : alert\_tier ENUM\\
        fused\_score : FLOAT\\
        reasons : TEXT[]\\
        explanation : JSONB\\
        supervisor\_decision : VARCHAR\\
        decision\_timestamp : TIMESTAMP
    };
    \end{tabular}
};

% ── evaluation_labels ──
\node[table, right=3cm of alerts] (eval) {
    \begin{tabular}{@{}c@{}}
    \tikz\node[header]{evaluation\_labels};\\[-0.4pt]
    \tikz\node[cols]{%
        \underline{event\_id} : UUID FK\\
        is\_anomaly : BOOLEAN\\
        anomaly\_type : VARCHAR
    };
    \end{tabular}
};

% ── profile_snapshots ──
\node[table, below=1.5cm of alerts] (profiles) {
    \begin{tabular}{@{}c@{}}
    \tikz\node[header]{profile\_snapshots};\\[-0.4pt]
    \tikz\node[cols]{%
        \underline{client\_id} : UUID\\
        \underline{snapshot\_date} : DATE\\
        profile\_data : JSONB
    };
    \end{tabular}
};

% ── employee_profile_snapshots ──
\node[table, right=3cm of profiles] (emp_profiles) {
    \begin{tabular}{@{}c@{}}
    \tikz\node[header]{employee\_profiles};\\[-0.4pt]
    \tikz\node[cols]{%
        \underline{employee\_id} : VARCHAR\\
        \underline{snapshot\_date} : DATE\\
        profile\_data : JSONB
    };
    \end{tabular}
};

% ── archetype_baselines ──
\node[table, below=1cm of profiles] (baselines) {
    \begin{tabular}{@{}c@{}}
    \tikz\node[header]{archetype\_baselines};\\[-0.4pt]
    \tikz\node[cols]{%
        \underline{archetype} : VARCHAR PK\\
        baseline\_data : JSONB
    };
    \end{tabular}
};

% ── Relationships ──
\draw[fk] (accounts.north) -- node[left, font=\tiny, pos=0.5] {FK} (clients.south);
\draw[fk] (eval.north) -- node[right, font=\tiny, pos=0.5] {FK} (transactions.south);
\draw[nofk] (alerts.north east) -- node[left, font=\tiny, pos=0.4] {no FK} (transactions.south west);

\end{tikzpicture}
}%
\caption{Database schema --- entity relationship diagram. Solid arrows represent
         foreign key constraints; the dashed arrow between \texttt{alerts} and
         \texttt{transactions} indicates a logical link without a database-level
         constraint (Section~\ref{sec:sinks}).}
\label{fig:er-diagram}
\end{figure}

\subsection{Master Tables}

Two master tables store the static reference data loaded once from the
generator output:

\begin{table}[H]
\centering
\caption{\texttt{clients\_master} --- one row per client}
\label{tab:clients-master}
\begin{tabularx}{\textwidth}{@{}llX@{}}
\toprule
\textbf{Column} & \textbf{Type} & \textbf{Description} \\
\midrule
\texttt{client\_id}           & UUID PK      & Unique client identifier \\
\texttt{archetype}            & VARCHAR(20)  & One of: salaried, small\_business, student, retiree, big\_business \\
\texttt{account\_opening\_date} & DATE        & Used to compute \texttt{account\_age\_days} during enrichment \\
\texttt{nationality}          & VARCHAR(5)   & Default \texttt{TN} (Tunisia) \\
\texttt{client\_type}         & VARCHAR(30)  & Regulatory classification: personne\_physique or personne\_morale \\
\texttt{home\_branch\_id}     & VARCHAR(10)  & Primary branch for branch-match feature \\
\bottomrule
\end{tabularx}
\end{table}

\begin{table}[H]
\centering
\caption{\texttt{employees\_master} --- one row per teller}
\label{tab:employees-master}
\begin{tabularx}{\textwidth}{@{}llX@{}}
\toprule
\textbf{Column} & \textbf{Type} & \textbf{Description} \\
\midrule
\texttt{employee\_id} & VARCHAR(10) PK & Teller identifier (e.g., \texttt{EMP-0001}) \\
\texttt{branch\_id}   & VARCHAR(10)    & Assigned branch (exactly 2 employees per branch) \\
\bottomrule
\end{tabularx}
\end{table}

\subsection{Transaction and Event Tables}

\begin{table}[H]
\centering
\caption{\texttt{transactions} --- populated by the raw events sink}
\label{tab:transactions-table}
\begin{tabularx}{\textwidth}{@{}llX@{}}
\toprule
\textbf{Column} & \textbf{Type} & \textbf{Description} \\
\midrule
\texttt{event\_id}       & UUID PK      & Unique event identifier \\
\texttt{client\_id}      & UUID         & Client performing the operation \\
\texttt{account\_id}     & UUID         & Target account \\
\texttt{employee\_id}    & VARCHAR(10)  & Teller processing the operation \\
\texttt{branch\_id}      & VARCHAR(10)  & Branch where the operation occurred \\
\texttt{timestamp}       & TIMESTAMP    & Operation timestamp \\
\texttt{amount}          & NUMERIC      & Transaction amount in TND \\
\texttt{currency}        & VARCHAR(3)   & Always \texttt{TND} \\
\texttt{channel}         & VARCHAR(20)  & Always \texttt{guichet} (teller window) \\
\texttt{operation\_type} & VARCHAR(20)  & One of: retrait, versement, virement, cheque \\
\texttt{payload}         & JSONB        & Operation-specific fields (beneficiary, cheque number, etc.) \\
\bottomrule
\end{tabularx}
\end{table}

\begin{table}[H]
\centering
\caption{\texttt{evaluation\_labels} --- ground truth for model evaluation}
\label{tab:eval-labels}
\begin{tabularx}{\textwidth}{@{}llX@{}}
\toprule
\textbf{Column} & \textbf{Type} & \textbf{Description} \\
\midrule
\texttt{event\_id}     & UUID FK $\rightarrow$ \texttt{transactions} & Links to the original event \\
\texttt{is\_anomaly}   & BOOLEAN  & Generator oracle label \\
\texttt{anomaly\_type} & VARCHAR  & Scenario name (e.g., smurfing, amount\_spike) or \texttt{NULL} \\
\bottomrule
\end{tabularx}
\end{table}

The \texttt{evaluation\_labels} table is separated from \texttt{transactions}
by design. Oracle labels are generator metadata used exclusively for model
evaluation --- they do not exist in a production system with real data. Keeping
them in a separate table makes it trivial to exclude them from production
deployments without modifying the transaction schema.

\subsection{Alerts Table}

\begin{table}[H]
\centering
\caption{\texttt{alerts} --- populated by the decisions sink}
\label{tab:alerts-table}
\begin{tabularx}{\textwidth}{@{}llX@{}}
\toprule
\textbf{Column} & \textbf{Type} & \textbf{Description} \\
\midrule
\texttt{alert\_id}            & SERIAL PK    & Auto-incrementing alert identifier \\
\texttt{event\_id}            & UUID         & References the triggering event (no FK --- see Section~\ref{sec:sinks}) \\
\texttt{client\_id}           & UUID         & Client involved \\
\texttt{employee\_id}         & VARCHAR(10)  & Teller who processed the operation \\
\texttt{branch\_id}           & VARCHAR(10)  & Branch location \\
\texttt{timestamp}            & TIMESTAMP    & Event timestamp \\
\texttt{amount}               & NUMERIC      & Transaction amount \\
\texttt{operation\_type}      & VARCHAR(20)  & Operation type \\
\texttt{tier}                 & alert\_tier ENUM & INFO, REVIEW, ALERT, or BLOCK \\
\texttt{fused\_score}         & FLOAT        & Combined AE + LSTM anomaly score \\
\texttt{reasons}              & TEXT[]       & Array of triggered rule names and score-based reasons \\
\texttt{regulatory\_flags}    & TEXT[]       & Regulatory rules that fired (logged regardless of tier impact) \\
\texttt{explanation}          & JSONB        & SHAP explanation with \texttt{reasons} (NLG text) and \texttt{raw} (feature-level detail) \\
\texttt{supervisor\_decision} & VARCHAR(10)  & \texttt{pending}, \texttt{confirmed}, or \texttt{rejected} \\
\texttt{supervisor\_notes}    & TEXT         & Free-text supervisor commentary \\
\texttt{decision\_timestamp}  & TIMESTAMP    & When the supervisor acted \\
\bottomrule
\end{tabularx}
\end{table}

The \texttt{tier} column uses a PostgreSQL ENUM type rather than a VARCHAR.
This enforces valid values at the database level and enables direct ordering
(\texttt{INFO < REVIEW < ALERT < BLOCK}) in queries without a CASE expression.

\subsection{Profile and Baseline Tables}

Three tables support the batch profile rebuild pipeline:

\begin{table}[H]
\centering
\caption{Profile and baseline tables}
\label{tab:profile-tables}
\begin{tabularx}{\textwidth}{@{}llX@{}}
\toprule
\textbf{Table} & \textbf{Key} & \textbf{Content} \\
\midrule
\texttt{profile\_snapshots}          & (\texttt{client\_id}, \texttt{snapshot\_date})
    & Client profile as JSONB: transaction statistics, behavioral distributions,
      financial state, risk history, recent events buffer \\
\addlinespace
\texttt{employee\_profile\_snapshots} & (\texttt{employee\_id}, \texttt{snapshot\_date})
    & Employee profile as JSONB: processed transaction stats, client distribution,
      peer comparison z-scores \\
\addlinespace
\texttt{archetype\_baselines}        & \texttt{archetype} PK
    & Per-archetype reference statistics as JSONB: mean amounts, operation mix,
      frequency norms. Used as the final fallback in the amount statistics
      fallback chain \\
\bottomrule
\end{tabularx}
\end{table}

Profiles are stored as JSONB rather than as individual columns. A client
profile contains approximately 194 fields across 6 categories; normalizing
these into columns would produce an unwieldy table that changes shape every
time a feature is added. JSONB stores the profile as a single document,
and the nightly rebuild overwrites it atomically via \texttt{ON CONFLICT}
upsert. The tradeoff is that individual fields cannot be indexed or queried
efficiently, but the access pattern (load entire profile into Redis) does not
require per-field queries.

% ═══════════════════════════════════════════════════════════
\section{Batch Integration}
\label{sec:batch-integration}
% ═══════════════════════════════════════════════════════════

The streaming pipeline depends on client profiles that are too expensive to
recompute on every event. A nightly batch pipeline rebuilds these profiles
from the full transaction history in PostgreSQL and loads the results into
Redis for the next day's streaming consumption.

\subsection{Nightly Profile Rebuild}

An Airflow DAG (\texttt{nightly\_profiles}) runs eight tasks with validation
gates between stages:

\begin{enumerate}
    \item Rebuild client profile snapshots from the \texttt{transactions} table
    \item Validate snapshot completeness (row counts, null checks)
    \item Rebuild employee profile snapshots
    \item Validate employee snapshots
    \item Recompute archetype baselines
    \item Validate baselines
    \item Hydrate Redis with updated profiles
    \item Run a smoke test (score a sample event with the new profiles)
\end{enumerate}

Each task uses \texttt{ON CONFLICT} upserts on its target table, making every
task individually retriable. If an Airflow task fails mid-execution and is
retried, the upsert ensures that partially-written rows are overwritten
rather than duplicated.

\subsection{Feature Parity Guarantee}

The batch profile computation calls the same
\texttt{EnrichmentCore.\_compute\_features()} method used by the streaming
pipeline. The method is a pure function: it takes an event, a profile, and
buffer data as arguments and returns 30 features with no side effects. The
batch script passes the archetype baseline as a parameter rather than
fetching it from Redis, which is the only difference from the streaming path.

This design guarantees that the features used to rebuild profiles overnight are
computed with identical logic to the features produced in real time. Any change
to the feature computation --- adding a feature, modifying a z-score formula,
adjusting the fallback chain --- is automatically reflected in both paths
because both paths call the same function.

\subsection{Streaming--Batch Coordination}

The current deployment does not pause streaming consumers during the batch
rebuild window. In practice, the nightly DAG runs during low-traffic hours
when transaction volume is minimal, and the Redis hydration step atomically
overwrites profile keys. A brief window exists where a streaming event could
read a stale profile (pre-rebuild) and then the next event reads the updated
profile, creating a one-event inconsistency. For the current deployment this
is acceptable; a production system would implement a coordination mechanism
such as pausing consumers during the hydration step or using Redis
transactions to swap profile versions atomically. This is documented as
future work in Chapter~\ref{ch:roadmap}.